% ============================================================================
%  Relatorio Tecnico - Trabalho 2 de Programacao Paralela (MPI coord/trab)
%  Formato exigido:
%    pag. 1  -> coluna DUPLA, margens 1cm, fonte 10, somente texto
%    pag. 2  -> coluna SIMPLES: tabelas + graficos (anexo)
%    pag. 3+ -> coluna SIMPLES: codigo fonte
%  Compilar a partir de relatorio/ (usa ../resultados/ e ../ para os fontes):
%    tectonic relatorio.tex    ou    pdflatex relatorio.tex
%  \pend{...} marca em VERMELHO o que depende das medicoes no cluster.
%  Antes de entregar: grep pend relatorio.tex deve voltar so a definicao.
%  Os numeros saem de resultados/tabelas.txt (python3 analisar.py).
% ============================================================================
\documentclass[10pt,a4paper,twocolumn]{article}

\usepackage{iftex}
\ifPDFTeX
  \usepackage[utf8]{inputenc}
  \usepackage[T1]{fontenc}
\fi
\usepackage[brazil]{babel}
\usepackage[margin=1cm]{geometry}
\usepackage{graphicx}
\usepackage{booktabs}
\usepackage{listings}
\usepackage{xcolor}
\usepackage{catchfile}

\setlength{\parindent}{0pt}
\setlength{\parskip}{3pt}
\setlength{\columnsep}{6mm}
\pagestyle{empty}

\newcommand{\pend}[1]{\textcolor{red}{\textbf{[#1]}}}

% Inclui o grafico se ja existir; senao deixa um aviso no lugar.
\newcommand{\grafico}[2][\linewidth]{\IfFileExists{../resultados/#2.pdf}%
  {\includegraphics[width=#1]{../resultados/#2.pdf}}%
  {\pend{falta \detokenize{#2}.pdf: rodar python3 graficos.py}}}

% Linhas de tabela geradas por analisar.py (resultados/tab_*.tex); enquanto
% nao houver medicao, mostra as linhas pendentes passadas como 2o argumento.
% Carregadas em macros ANTES das tabelas: um \input direto dentro de um
% tabular quebra o \bottomrule ("Misplaced \noalign").
\newcommand{\carrega}[1]{\IfFileExists{../resultados/#1.tex}%
  {\expandafter\CatchFileDef\csname #1\endcsname{../resultados/#1.tex}{}}{}}
\carrega{tab_forte}\carrega{tab_fraca}\carrega{tab_decomposicao}\carrega{tab_grao}
\newcommand{\linhas}[2]{\ifcsname #1\endcsname\csname #1\endcsname\else #2\fi}

\lstdefinestyle{codigoC}{
  language=C, basicstyle=\ttfamily\scriptsize, numbers=left,
  numberstyle=\tiny\color{gray}, keywordstyle=\color{blue!70!black},
  commentstyle=\itshape\color{green!40!black}, stringstyle=\color{red!60!black},
  showstringspaces=false, breaklines=true, tabsize=4,
  columns=fullflexible, lineskip=-1pt,
  morekeywords={MPI_Init,MPI_Send,MPI_Recv,MPI_Status,MPI_Gather,MPI_Gatherv,
                MPI_Barrier,MPI_Wtime}
}

\begin{document}

% ============================ CABECALHO =====================================
\twocolumn[
{\centering
  {\large\bfseries Mandelbrot com MPI no Modelo Coordenador/Trabalhador}\\[2pt]
  {\small Trabalho 2 --- Programação Paralela --- Grupo 20 (conta LAD \texttt{cp32020})}\\[1pt]
  {\small Gabriel Coelho \quad Gabriel Cherubini}\\[1pt]
  {\small Escola Politécnica --- PUCRS \quad --- \quad 2026/2}\par}
\vspace{3mm}
]

% ======================== PAGINA 1: TEXTO ===================================

\textbf{1. Problema e mapeamento.}
Mantivemos o problema do T1: para cada pixel $c$ de uma janela fixa do plano
complexo, itera-se $z_{n+1}=z_n^2+c$ até $|z|>2$ ou 2000 iterações. O
\textbf{saco de trabalho} é a imagem, uma matriz de inteiros que só o
coordenador guarda; cada \textbf{vetor} é uma linha, e processá-lo é calcular
as iterações dos seus pixels. As linhas são independentes, mas de custo muito
desigual (Fig.~3, esq.): as centrais cruzam o conjunto e custam o máximo por
pixel, as das bordas escapam logo. É o cenário em que distribuir sob demanda
compensa.

\textbf{2. Implementação SPMD.}
Um fonte gera as duas versões (\texttt{gcc} e \texttt{mpicc -DPARALELO}, ambas
\texttt{-O2}). O rank 0 coordena e não calcula; os demais trabalham. O
trabalhador toma a iniciativa com \texttt{PEDIDO}; recebe \texttt{TAREFA} com
o índice da linha; devolve \texttt{RESULTADO} = [índice, pixels\ldots], que já
vale como novo pedido. O coordenador recebe com \texttt{MPI\_ANY\_SOURCE}
(\textbf{ordem de chegada}) e usa o índice para gravar o vetor na sua linha,
\textbf{mantendo a ordem do saco}; vazio o saco, responde \texttt{FIM}.
Funciona para qualquer $p\ge2$; \texttt{-k} agrupa $k$ linhas por tarefa. O
checksum (soma das iterações) é \textbf{idêntico} ao sequencial em todas as
execuções \pend{checksum}. A instrumentação mede, em cada processo, o tempo
calculando, bloqueado em \texttt{Recv}/\texttt{Send} e o instante de
dispensa, e registra o início e fim de cada tarefa; é coletada após o tempo
medido, sem afetá-lo.

\textbf{3. Ambiente e metodologia.}
Dois nós exclusivos da atlantica \pend{nós}, cada um com 2$\times$ Xeon
\pend{modelo} (8 núcleos, 16 threads HT), \pend{MPI}, gcc \pend{versão}.
$p=2..32$ processos alternados entre os nós; até 16 há um por núcleo físico,
com 32 usam-se as HT. Cada tempo é a mediana de 5 execuções (barras min--máx
na Fig.~1); $T_{seq}$ é o binário sequencial. Como $p$ inclui o coordenador
ocioso, o teto do speed-up é $p-1$ e o da eficiência, $(p-1)/p$.

\textbf{4. Escalabilidade forte e fraca.}
\pend{S(16)=?, E=?\%; S(32)=?, E=?\%, contra teto (p-1)/p. Fraca: E(32)=?\%.
Uma frase ligando à seção 5. Comparar com o T1 em 1 nó: OpenMP S(8)=6,90.}

\textbf{5. Para onde vai a eficiência.}
Dividimos os recursos alocados, $p\cdot T$, em parcelas medidas (Fig.~2,
Tab.~3): trabalho útil ($T_{seq}$, que é a própria eficiência); cálculo mais
lento que no sequencial; o coordenador; comunicação; ociosidade final. A soma
fecha em 100\% (resíduo < \pend{?}\%).
\pend{Ler os números: até p=16 a perda é quase só o coordenador (1/p);
comunicação ?\%; com 32 surge ?\% de cálculo mais lento = duas threads HT
dividindo o núcleo. O coordenador fica ?\% ocioso: não é gargalo.}

\textbf{6. Balanceamento de carga.}
\pend{Com k=1, tempo calculando máx/média = ? (CV ?\%), embora as linhas por
trabalhador variem de ? a ?: quem pega linhas caras recebe menos. Com um bloco
fixo por trabalhador (k=133, divisão estática) máx/média = ? e o tempo sobe
?$\times$. O perfil de custo previa ? s; medido ? s: o desbalanceamento vem do
custo das linhas. O Gantt (Fig.~4) mostra todos terminando juntos no saco e a
cauda no estático. Grão (Tab.~4): k=? é o melhor; k grande perde balanceamento.}

\textbf{7. Comunicação e rede.}
\pend{Intensidade (Fig.~5 esq.): com ? µs de cálculo por mensagem E=?\%;
acima de ? µs se aproxima do teto. 1×2 nós: p=8 em 1 nó ?\% × 2 nós ?\%
(custo da rede); p=16 em 1 nó com HT ?\% × 2 nós ?\%: HT custa mais/menos
que a rede.}

\textbf{8. Uso de IA.}
Claude Code (Anthropic): implementação do protocolo, instrumentação, scripts
de medição e análise, gráficos e esqueleto do relatório. O grupo definiu o
mapeamento e a metodologia, executou e validou as medições e redigiu e
revisou a análise.

\textbf{9. Conclusão.}
\pend{Resumo em 3 frases com os números principais.}

% ==================== PAGINA 2: TABELAS E GRAFICOS ==========================
\onecolumn
\newpage
\small

\begin{minipage}[t]{0.49\linewidth}
\textbf{Tab. 1 --- Forte.} 4096$^2$, 2000 it, $k=1$, mediana de 5.
$S=T_{seq}/T$, $E=S/p$.\\[2pt]
\begin{tabular}{rrrrr}
\toprule
$p$ & Tempo (s) & $S$ & $E$ (\%) & teto (\%) \\
\midrule
\linhas{tab_forte}{%
seq & \pend{?} & --- & --- & --- \\
2  & \pend{?} & \pend{?} & \pend{?} & 50,0 \\
4  & \pend{?} & \pend{?} & \pend{?} & 75,0 \\
8  & \pend{?} & \pend{?} & \pend{?} & 87,5 \\
16 & \pend{?} & \pend{?} & \pend{?} & 93,8 \\
32 & \pend{?} & \pend{?} & \pend{?} & 96,9 \\}
\bottomrule
\end{tabular}
\end{minipage}\hfill
\begin{minipage}[t]{0.49\linewidth}
\textbf{Tab. 2 --- Fraca.} 1024$^2$ pixels por processo.
$S=p\,T_{seq}/T$, $E=T_{seq}/T$.\\[2pt]
\begin{tabular}{rlrrr}
\toprule
$p$ & Imagem & Tempo (s) & $S$ & $E$ (\%) \\
\midrule
\linhas{tab_fraca}{%
seq & 1024$^2$ & \pend{?} & 1,00 & 100,0 \\
2  & 1448$^2$ & \pend{?} & \pend{?} & \pend{?} \\
4  & 2048$^2$ & \pend{?} & \pend{?} & \pend{?} \\
8  & 2896$^2$ & \pend{?} & \pend{?} & \pend{?} \\
16 & 4096$^2$ & \pend{?} & \pend{?} & \pend{?} \\
32 & 5793$^2$ & \pend{?} & \pend{?} & \pend{?} \\}
\bottomrule
\end{tabular}
\end{minipage}

\vspace{3mm}
\begin{minipage}[t]{0.49\linewidth}
\textbf{Tab. 3 --- Decomposição de $p\cdot T$} (\%), $k=1$.\\[2pt]
\begin{tabular}{rrrrrr}
\toprule
$p$ & útil & cálc.+lento & coord. & comun. & ocioso fim \\
\midrule
\linhas{tab_decomposicao}{%
2  & \pend{?} & \pend{?} & \pend{?} & \pend{?} & \pend{?} \\
4  & \pend{?} & \pend{?} & \pend{?} & \pend{?} & \pend{?} \\
8  & \pend{?} & \pend{?} & \pend{?} & \pend{?} & \pend{?} \\
16 & \pend{?} & \pend{?} & \pend{?} & \pend{?} & \pend{?} \\
32 & \pend{?} & \pend{?} & \pend{?} & \pend{?} & \pend{?} \\}
\bottomrule
\end{tabular}
\end{minipage}\hfill
\begin{minipage}[t]{0.49\linewidth}
\textbf{Tab. 4 --- Grão}, $p=32$. $k=133$: um bloco fixo por trabalhador
(estático).\\[2pt]
\begin{tabular}{rrrrr}
\toprule
$k$ & Tempo (s) & $E$ (\%) & ocup. máx/méd & vs melhor \\
\midrule
\linhas{tab_grao}{%
1   & \pend{?} & \pend{?} & \pend{?} & \pend{?} \\
4   & \pend{?} & \pend{?} & \pend{?} & \pend{?} \\
16  & \pend{?} & \pend{?} & \pend{?} & \pend{?} \\
64  & \pend{?} & \pend{?} & \pend{?} & \pend{?} \\
133 & \pend{?} & \pend{?} & \pend{?} & \pend{?} \\}
\bottomrule
\end{tabular}
\end{minipage}

\vspace{3mm}
\begin{center}
\grafico[0.92\linewidth]{grafico_escala}\\
\textbf{Fig. 1} --- Speed-up (barras: mín--máx de 5 execuções) e eficiência.
Tracejado: ideal $S=p$; pontilhado: teto do modelo, $p-1$.\\[2mm]
\grafico[0.92\linewidth]{grafico_decomposicao}\\
\textbf{Fig. 2} --- Recursos alocados ($p\times T$) divididos pelo que
cada processo fez; a fatia azul é a eficiência.
\end{center}

\newpage
\begin{center}
\grafico[0.92\linewidth]{grafico_balanceamento}\\
\textbf{Fig. 3} --- Esq.: custo de cada linha, com as fronteiras dos blocos
estáticos. Dir.: tempo calculando de cada trabalhador, saco ($k=1$) ×
bloco fixo ($k=133$).\\[2mm]
\grafico[0.92\linewidth]{grafico_gantt}\\
\textbf{Fig. 4} --- Linha do tempo das tarefas de cada trabalhador ($p=32$).\\[2mm]
\grafico[0.92\linewidth]{grafico_comunicacao}\\
\textbf{Fig. 5} --- Esq.: eficiência conforme o cálculo por mensagem cresce
(mesmo número de mensagens, \texttt{max\_iter} variando). Dir.: o mesmo $p$
em 1 e em 2 nós.
\end{center}
\normalsize

% ==================== CODIGO FONTE ==========================================
\newpage
\begin{center}\textbf{\large Código fonte --- \texttt{mandelbrot.c}}\end{center}
\lstinputlisting[style=codigoC]{../mandelbrot.c}

\newpage
\begin{center}\textbf{\large Medição --- \texttt{medir.sh} e \texttt{job.sh}}\end{center}
\lstinputlisting[style=codigoC,language=bash]{../medir.sh}
\lstinputlisting[style=codigoC,language=bash]{../job.sh}

\newpage
\begin{center}\textbf{\large Análise --- \texttt{analisar.py}}\end{center}
\lstinputlisting[style=codigoC,language=Python]{../analisar.py}

\end{document}
